% problem-000147 6 同位素效应与动力学
\section*{6. 同位素效应与动力学}
\textit{下列为分问。}

\subsection*{6-1}
动⼒学同位素效应可⽤于反应机理研究。⼈们研究了苯基异丙基酮溴化反应的动⼒学，298 K的数据如下表，设其速率⽅程为 $r = k _ { \neq \neq \forall \downarrow } [ \mathrm { K } - \mathrm { H } ] ^ { \alpha } [ \mathrm { H } ^ { + } ] ^ { \beta } [ \mathrm { B r } _ { 2 } ] ^ { \gamma }$ ，指出各反应物的分级数并计算表观速率常数。

\textit{（原卷此处含表格/仪器试剂表，详见 PDF 或 MD 源文件。）}

如下图所⽰，K-H表⽰苯基异丙基酮分⼦，K-D表⽰其氘取代分⼦。

（图：）
K-H

（图：）
图 6-1
K-D

\subsection*{6-2}
苯基异丙基酮的溴化是酸催化反应。⼀般认为可能的机理为：

$
\begin{array}{l} \mathrm {K - H+ H^ {+} \xleftrightarrow [ k_{-1} ]{k_ {1}} HK- H^ {+}} \\ \mathrm {HK- H^ {+} \xrightarrow {k_ {2}} Enol(烯醇)+ H^ {+}} \\ \mathrm {Enol + Br_ {2} \xrightarrow {k_ {3}} K - Br+ H^ {+} + Br^ {-}} \end{array}
$

采⽤稳态近似，做合理近似，推导总反应的速率⽅程。指出此反应的速控步并说明理由。

\subsection*{6-3}
保持反应物浓度以及缓冲溶液pH值不变，308 K下苯基异丙基酮溴化反应速率相⽐298 K时加倍，试估算该反应速控步的活化能。已知反应 $\mathrm { K \mathrm { - } H + H ^ { + }  H K } .$ -H+的Δ� = 30 kJ mol−1。

\subsection*{6-4}
以C—H键解离反应的动⼒学同位素效应为例，说明如下。键伸缩振动可近似⽤简谐振动描述，简谐振动频率 $\begin{array} { r } { \mathsf { v } = \frac { 1 } { 2 \pi } \sqrt { \frac { \kappa } { \mu } } , } \end{array}$ 式中�为⼒常数（C—H与C—D的⼒常数�相等），�为折合质量，以C—H为例： $\mu =$ $\frac { m _ { \mathrm { C } } m _ { \mathrm { H } } } { m _ { \mathrm { C } } + m _ { \mathrm { H } } } { \mathrm { ^ c } }$ 。分⼦在振动基态时的能量称为零点能 $E _ { 0 }$ ， $\begin{array} { r } { E _ { 0 } = \frac 1 2 h \nu } \end{array}$ ，式中ℎ为普朗克常数。如下图所⽰，C—D与C—H反应物的零点能不同，但是因过渡态伸缩振动的⼒常数很⼩，C—D与C—H过渡态(TS)零点能的差值通常可以忽略，因此，C—D键解离⽐C—H键解离的活化能更⼤。近似认为两种反应的Arrhenius指前因⼦相等，则反应速率常数的变化可以定量计算。计算K—H和K—D在298K时溴化反应速控步的速率常数⽐值。已知$\widetilde { \nu } _ { \mathrm { C - H } } = 2 9 5 0 \mathrm { c m } ^ { - 1 }$ ，H，D相对原⼦质量分别为1.007，2.000。

（图：）
图 6-2

\subsection*{6-5}
碱性条件下，K-H氯代反应的反应机理会发⽣变化，总反应的表观速率⽅程为： $r = k [ \mathrm { C l ^ { - } } ] [ \mathrm { O H ^ { - } } ] [ \mathrm { K \mathrm { - } H } ]$ 试拟定K-H氯代反应的反应机理。

\subsection*{6-6}
氘代对化学反应平衡常数的影响也可以从零点能出发，类似讨论。已知298K下，醋酸电离反应HOAc$+ \mathrm { H } _ { 2 } \mathrm { O }  \mathrm { O A c } ^ { - } + \mathrm { H } _ { 3 } \mathrm { O } ^ { + }$ 的 $\mathsf { p } K _ { a } = 4 . 7 4$ $\mathrm { H _ { 3 } O ^ { + } }$ 与HOAc中O—H键的振动波数分别为 $3 0 0 0 \mathrm { c m } ^ { - 1 }$ 和 $2 8 0 0 \mathrm { c m } ^ { - }$ $^ 1 \circ$ 判断 HOAc 在 $\mathrm { D } _ { 3 } \mathrm { O } ^ { + }$ 中电离反应 $\mathsf { p } K _ { a }$ 相⽐ HOAc 在 $\mathrm { H _ { 3 } O ^ { - } }$ +中电离反应的 $\mathsf { p } K _ { a }$ 如何变化并说明理由（不要求计算$\mathsf { p } K _ { a }$ 数值）。
